\section{Determinación de propiedades opticas de eritrocitos}
\label{section:flow_cytometry}

La citometría de flujo es una técnica experimental que permite analizar y diferenciar grandes poblaciones celulares (tasa de $10^{6}$ células por minuto para muestras diluidas \cite{ricoFlowCytometry2021}), particularmente eritrocitos, a partir de la luz esparcida por las células al atravesar un láser \cite{bectondickinsonIntroductionFlowCytometry2000}. La información obtenida mediante esta técnica proviene de la distribución angular de la luz esparcida por cada célula, la cual depende de propiedades como su volumen, forma, orientación e índice de refracción \cite{bectondickinsonIntroductionFlowCytometry2000}. En los citómetros convencionales se registran dos señales de esparcimiento: el esparcimiento frontal (FS, por sus siglas en inglés), medido en un intervalo angular pequeño (menor a $20^{\circ}$ \cite{ricoFlowCytometry2021}) alrededor de la dirección de propagación del haz incidente ($\theta=0^{\circ}$) y asociado principalmente con el tamaño celular, y el esparcimiento lateral (SS, por sus siglas en inglés), medido usualmente a $\theta=90^\circ$  respecto al haz incidente \cite{ricoFlowCytometry2021} y relacionado con la complejidad interna o granularidad de la célula \cite{mckinnonFlowCytometryOverview2018}. 

En la Fig.~\ref{fig:flow_cyt} se muestra un arreglo experimental característico empleado en citometría de flujo. Una suspensión de células es transportada mediante un sistema de microfluídica~\cite{kageFlowCytometryMicrofluidic2018} que permite controlar su flujo y favorecer que atraviesen individualmente la región de análisis~\cite{mckinnonFlowCytometryOverview2018}. En esta región, el flujo celular se ilumina mediante un láser polarizado \cite{asburyPolarizationScatterFluorescence2000}, de manera que la luz esparcida por cada célula pueda registrarse de forma individual. En años recientes, se han desarrollado citómetros con enfoque acústico, que permiten mejorar la alineación y orientación de las células durante su paso por la región de análisis~\cite{giengerAssessmentDeformationHuman2019}. Esta capacidad resulta especialmente relevante en el estudio de eritrocitos, ya que su respuesta de esparcimiento depende de la orientación de la célula respecto al campo incidente. La luz esparcida por las células se colecta mediante lentes convergentes y se dirige hacia detectores, como fotodiodos o tubos fotomultiplicadores \cite{mckinnonFlowCytometryOverview2018}. En la detección de FS, el haz incidente directo debe bloquearse para evitar que alcance el detector y lo sature \cite{arkesteijnImprovedFlowCytometric2020}. Para ello, se emplea una barra de bloqueo de luz (\textit{obscuration bar} o \textit{beam stop}) colocada en el sistema óptico de detección frontal, que bloquea la contribución directa del láser, mientras que la luz esparcida por la célula a ángulos pequeños puede pasar alrededor de la barra y ser recolectada por la lente hacia el detector de FS \cite{arkesteijnImprovedFlowCytometric2020}. Para bloquear la luz que llega con otros ángulos (esparcimiento de otras partículas ajenas a las de interés en la suspensión) se coloca un diafragma \cite{arkesteijnImprovedFlowCytometric2020}. Para la señal SS, la contribución directa del láser no incide sobre el detector, por lo que únicamente se coloca un diafragma para que la luz esparcida pueda detectarse \cite{arkesteijnImprovedFlowCytometric2020}. %

El arreglo experimental de los citómetros convencionales implica que registren únicamente ciertas regiones de la distribución angular del esparcimiento. Por ello, resulta de interés estudiar intervalos distintos de los asociados a FS y SS, con el fin de identificar ventanas angulares en las que las diferencias entre eritrocitos sanos y enfermos sean mayores. Con este objetivo, requiere la implementación de modelos que contemplen con fidelidad la morfología y la composición interna de los eritrocitos sanos y enfermos.


\begin{figure}[tb]
	\centering
	\includegraphics[width=9cm]{2-Eritrocitos/Flow_cytometry.pdf}
	
	\caption{Esquema de un citómetro de flujo convencional. Las células a analizar viajan en una suspensión hasta intersectar con un láser polarizado. Para la medición del FS, se coloca una barra de bloqueo de luz para bloquear el haz incidente y posteriormente, se colecta por una lente convergente y pasa por un diafragma para evitar que el detector capture la luz proveniente de otros ángulos, asociada al esparcimiento de otras partículas ajenas a las células de interés. Para la medición del SS la luz se colecta por una lente convergente, pasa por un diafragma y se mide por medio de un detector.} 
	\label{fig:flow_cyt}
\end{figure}
%